\documentclass[11pt,a4paper]{article}

% ----------------------------- Encodage / langue -----------------------------
\usepackage[utf8]{inputenc}
\usepackage{lmodern}
\usepackage[T1]{fontenc}
\usepackage[french]{babel}
\usepackage[protrusion=true,expansion=false]{microtype}

% ----------------------------- Maths -----------------------------
\usepackage{amsmath,amssymb,amsthm,mathtools}

% ----------------------------- Mise en page -----------------------------
\usepackage[a4paper,margin=2.4cm]{geometry}
\usepackage{graphicx}
\usepackage{xcolor}
\usepackage{booktabs}
\usepackage{enumitem}
\usepackage{fancyhdr}
\usepackage{caption}
\usepackage[colorlinks=true,linkcolor=NavyBlue,citecolor=NavyBlue,urlcolor=NavyBlue]{hyperref}

\definecolor{NavyBlue}{HTML}{1F5FA8}
\definecolor{Accent}{HTML}{C0392B}
\definecolor{LightGray}{HTML}{F4F6F9}
\definecolor{CodeKw}{HTML}{1F5FA8}
\definecolor{CodeStr}{HTML}{218C5A}
\definecolor{CodeCom}{HTML}{888888}

% ----------------------------- Code -----------------------------
\usepackage{listings}
\lstset{
  language=Python,
  basicstyle=\ttfamily\footnotesize,
  keywordstyle=\color{CodeKw}\bfseries,
  stringstyle=\color{CodeStr},
  commentstyle=\color{CodeCom}\itshape,
  numbers=left, numberstyle=\tiny\color{gray}, numbersep=8pt,
  backgroundcolor=\color{LightGray},
  frame=single, rulecolor=\color{gray!40},
  breaklines=true, showstringspaces=false,
  xleftmargin=1.6em, framexleftmargin=1.6em,
  literate={é}{{\'e}}1 {è}{{\`e}}1 {à}{{\`a}}1 {ê}{{\^e}}1 {î}{{\^i}}1 {ô}{{\^o}}1 {ç}{{\c c}}1 {û}{{\^u}}1 {É}{{\'E}}1
}

% ----------------------------- Théorèmes -----------------------------
\theoremstyle{plain}
\newtheorem{proposition}{Proposition}
\newtheorem{lemme}{Lemme}
\newtheorem{theoreme}{Théorème}
\theoremstyle{remark}
\newtheorem*{remarque}{Remarque}

% ----------------------------- En-têtes -----------------------------
\pagestyle{fancy}
\fancyhf{}
\renewcommand{\headrulewidth}{0.4pt}
\fancyhead[L]{\small\itshape Résolution d'une EDO d'ordre 2}
\fancyhead[R]{\small ING1 — 2025--2026}
\fancyfoot[C]{\thepage}

% ----------------------------- Macros -----------------------------
\newcommand{\R}{\mathbb{R}}
\newcommand{\eps}{\varepsilon}
\newcommand{\dd}{\,\mathrm{d}}

\begin{document}

% =========================================================================
% PAGE DE TITRE
% =========================================================================
\begin{titlepage}
{\raggedleft
\includegraphics[width=3.2cm]{logo_cytech.png}\par}
\centering
\vspace*{0.4cm}
{\large CY Tech — CY Cergy Paris Université}\\[0.3cm]
{\normalsize Filière Mathématiques Appliquées et Science des Données}\\[0.2cm]
{\normalsize Année universitaire 2025--2026}\\[2.2cm]

\rule{\linewidth}{0.4pt}\\[0.5cm]
{\Huge\bfseries Résolution d'une équation différentielle d'ordre deux}\\[0.6cm]
{\Large Analyse théorique, synthèse et implémentation numérique}\\[0.4cm]
\rule{\linewidth}{0.4pt}\\[0.3cm]
\vspace{2.0cm}

{\large\bfseries Étudiants}\\[0.4cm]
\begin{tabular}{c}
  Arun Kuganesan \quad\textbullet\quad Haytham Bouhou \quad\textbullet\quad Jérémie Konda\\[0.25cm]
  Amine Taki \quad\textbullet\quad Ulysse Beqiraj\\
\end{tabular}\\[2.0cm]

{\large\bfseries Projet GM \& MF — ING1}\\[2.2cm]

\vfill
{\normalsize Date de rendu : juin 2026}
\end{titlepage}

% =========================================================================
% RÉSUMÉ
% =========================================================================
\thispagestyle{fancy}
\section*{Résumé}
\addcontentsline{toc}{section}{Résumé}

Ce projet étudie une équation différentielle d'ordre deux sur la droite réelle,
dépendant de deux paramètres strictement positifs $\mu$ et $c$. La question
centrale est de déterminer pour quelles valeurs de ces paramètres il existe une
solution réalisant une \emph{transition} (un \emph{front}) reliant le palier $c$
en $-\infty$ au palier $0$ en $+\infty$.

Le travail comporte trois volets. Le premier, \textbf{analytique}, montre qu'une
telle solution impose des contraintes précises sur $(\mu,c)$, exprimées au moyen
d'une fonction auxiliaire $P$ : ceci justifie un algorithme de recherche des
paramètres. Le deuxième, de \textbf{synthèse}, établit la réciproque : sous ces
contraintes, une solution existe et est unique à translation près. Le troisième
propose une \textbf{implémentation Python} des deux fonctions demandées. Nous
exhibons en outre une \emph{solution analytique exacte} dans le cas modèle, qui
sert de référence pour valider le code (erreur de l'ordre de $10^{-10}$), et nous
discutons de façon critique le domaine de validité de l'algorithme. Le problème
se rattache à la théorie des fronts d'onde des équations de réaction–diffusion
de type \emph{Allen–Cahn / Nagumo}.

\newpage
\tableofcontents
\newpage

% =========================================================================
% 1. INTRODUCTION ET ÉTAT DE L'ART
% =========================================================================
\section{Introduction et état de l'art}

\subsection{Contexte}
L'équation étudiée,
\[
  -Q'' + \mu Q + g(Q) = 0 \quad \text{sur } \R,
  \qquad Q(-\infty)=c,\ Q(+\infty)=0,
\]
est l'\emph{équation de profil} associée aux ondes stationnaires d'une équation
d'évolution de type réaction–diffusion
\[
  u_t = u_{xx} - \mu u - g(u).
\]
En cherchant une onde progressive de vitesse nulle $u(t,x)=Q(x)$, on retombe
exactement sur notre problème : $Q$ décrit le \emph{front} qui relie deux états
d'équilibre. Ces équations modélisent des phénomènes de transition de phase, de
propagation en génétique des populations et de transmission de l'influx nerveux.

\subsection{Travaux de référence}
La structure mathématique sous-jacente — une non-linéarité \emph{bistable}
$R(u)=\mu u+g(u)$ admettant trois zéros — est celle de l'équation
d'\textbf{Allen–Cahn}, introduite en 1979 pour décrire le mouvement des parois
d'antiphase dans les alliages~\cite{allencahn1979}, et de l'équation de
\textbf{Nagumo} en neurophysiologie~\cite{nagumo1962}. L'analyse rigoureuse de
l'existence, de l'unicité (à translation près) et de la stabilité de ces fronts
remonte aux travaux fondateurs de \textbf{Fife et McLeod}~\cite{fifemcleod1977},
qui emploient une méthode de plan de phase, et de \textbf{Aronson et
Weinberger}~\cite{aronsonweinberger1978}. Historiquement, l'étude des fronts de
réaction–diffusion débute avec l'équation de
\textbf{Kolmogorov–Petrovsky–Piskunov}~\cite{kpp1937}, dans le cas monostable.

Le présent projet retrouve, par des moyens élémentaires (conservation d'une
\emph{énergie}, séparation des variables, théorème de Cauchy–Lipschitz), un cas
particulier de ces résultats : l'existence et l'unicité d'un front pour le
potentiel à \emph{puits de profondeur égale}. La condition d'égalité des puits
correspond à la condition dite \emph{de Maxwell}, classique en thermodynamique
des transitions de phase. Les outils analytiques mobilisés (méthode de l'énergie,
plan de phase) sont standard et exposés dans~\cite{evans2010,strogatz2015}.

\subsection{La condition de Maxwell}
\label{sec:maxwell}
La condition $P'(c)=0$ admet une interprétation thermodynamique éclairante.
Posons $R(y)=\mu y+g(y)$ le terme de réaction. La condition $\mu=P(c)$, soit
$f(c)=\mu c^2+2G(c)=0$, se réécrit
\[
  \int_0^c R(y)\dd y=\Big[\tfrac{\mu}{2}y^2+G(y)\Big]_0^c
  =\tfrac{\mu}{2}c^2+G(c)=0.
\]
Autrement dit, \emph{l'aire algébrique de $R$ entre les deux équilibres $0$ et
$c$ est nulle}. C'est la célèbre \emph{condition de Maxwell} (ou règle des aires
égales) : les deux « puits » du potentiel
$W(y)=-\int_0^y R$ ont la même profondeur. Physiquement, c'est exactement le cas
où aucune des deux phases n'est favorisée, ce qui produit un front
\emph{stationnaire} (vitesse nulle). Nous reviendrons sur le rôle de cette
symétrie au paragraphe~\ref{sec:nonsym}, en discutant ce qui se passe lorsqu'elle
est rompue.

\subsection*{Applications}
Au-delà des mathématiques, ces fronts modélisent des phénomènes physiques
concrets : la \emph{séparation de phases} dans les alliages métalliques
(Allen–Cahn), la \emph{propagation de l'influx nerveux} le long d'un axone
(Nagumo), la \emph{diffusion de gènes} dans une population (KPP), ou encore les
\emph{fronts de combustion}. Dans tous ces cas, $Q$ représente le profil
spatial qui sépare deux états stables, et la question « pour quels paramètres un
tel profil existe-t-il ? » est précisément celle que ce projet résout.

\subsection{Objectifs et plan}
Nos objectifs sont : (i) démontrer la nécessité des conditions de l'algorithme
(Partie~\ref{sec:analyse}) ; (ii) en établir la suffisance et l'unicité
(Partie~\ref{sec:synthese}) ; (iii) implémenter et valider numériquement la
méthode (Partie~\ref{sec:numerique}).

% =========================================================================
% 2. PRÉSENTATION DU PROBLÈME
% =========================================================================
\section{Présentation du problème}

\subsection{Cadre et notations}
On se donne $g\in\mathcal{C}^2(\R,\R)$ vérifiant $g(0)=g'(0)=0$, et l'on pose
\[
  G(y)=\int_0^y g(t)\dd t.
\]
On cherche une fonction $Q\in\mathcal{C}^2(\R)$ et des paramètres $\mu,c>0$ tels
que
\begin{equation}\label{eq:edo}
  \begin{cases}
    -Q''(x)+\mu Q(x)+g(Q(x))=0 & \forall x\in\R,\\[2pt]
    Q(+\infty)=0,\quad Q(-\infty)=c.
  \end{cases}
\end{equation}
Intuitivement, on cherche un \emph{profil de transition} régulier descendant de
$c$ vers $0$, avec deux paliers horizontaux aux extrémités.

\subsection{L'algorithme proposé}
On introduit $P:(0,+\infty)\to\R$ définie par
\[
  P(x)=-\frac{2G(x)}{x^2}.
\]
L'algorithme recherche les valeurs $c>0$ telles que
\begin{enumerate}[label=(\roman*),leftmargin=2.2em]
  \item $P'(c)=0$ ;
  \item $P(c)>0$ ;
  \item $P(x)<P(c)$ pour tout $x\in[0,c[$.
\end{enumerate}
Pour un tel $c$, on pose $\mu=P(c)$. L'objectif est de prouver que cet algorithme
capture \emph{exactement} toutes les solutions de~\eqref{eq:edo}.

% =========================================================================
% 3. PARTIE 1 : ANALYSE
% =========================================================================
\section{Partie 1 : Analyse}
\label{sec:analyse}

Dans toute cette partie, on suppose qu'il existe $Q\in\mathcal{C}^2(\R)$ solution
de~\eqref{eq:edo} pour des paramètres $\mu,c>0$.

\subsection{Une quantité conservée}
On définit
\[
  H(x)=\frac{Q'(x)^2}{2}-G(Q(x))-\frac{\mu}{2}Q(x)^2,
\]
qui joue le rôle d'une \emph{énergie mécanique} (terme cinétique plus terme
potentiel). En dérivant,
\[
  H'(x)=Q'(x)\big[Q''(x)-\mu Q(x)-g(Q(x))\big]=0
\]
d'après~\eqref{eq:edo}, donc $H$ est constante. En évaluant en $+\infty$, où
$Q\to0$ (donc $G(Q)\to G(0)=0$) et $Q'\to0$, on obtient $H\equiv0$, c'est-à-dire
l'identité centrale
\begin{equation}\label{eq:energie}
  Q'(x)^2 = \mu Q(x)^2 + 2G(Q(x)) =: f(Q(x)),
  \qquad f(y):=\mu y^2+2G(y).
\end{equation}

\begin{remarque}
La convergence $Q'\to0$ se justifie ainsi : comme $Q''=\mu Q+g(Q)$ admet une
limite finie en $+\infty$, si cette limite était non nulle alors $Q'$, et donc
$Q$, divergerait — ce qui contredit $Q\to0$. La limite de $Q''$ est donc nulle,
et celle de $Q'$ aussi.
\end{remarque}

\subsection{Positivité du second membre près de $0$}
Comme $g(0)=g'(0)=0$, la formule de Taylor donne $g(t)=O(t^2)$, d'où
$G(y)=O(y^3)$. Ainsi, pour $y>0$,
\[
  f(y)=y^2\!\left(\mu+\frac{2G(y)}{y^2}\right)\xrightarrow[y\to0^+]{}\ \mu\,y^2>0.
\]
Le second membre de~\eqref{eq:energie} est donc strictement positif pour $y>0$
assez petit.

\subsection{Réduction à une EDO d'ordre 1}
De~\eqref{eq:energie} on tire $Q'(x)=\eps(x)\sqrt{f(Q(x))}$ avec
$\eps(x)\in\{-1,+1\}$. Montrons que $\eps$ est \emph{globalement constant}.
Comme $Q'$ est continue, un changement de signe de $\eps$ imposerait l'annulation
de $Q'$ en un point $x_0$, donc $f(Q(x_0))=0$ ; le couple $(Q(x_0),0)$ serait une
donnée de Cauchy de l'EDO d'ordre 2. \emph{Si} $Q(x_0)$ est un point d'équilibre
(c'est-à-dire $\mu Q(x_0)+g(Q(x_0))=0$), la solution constante convient et
Cauchy–Lipschitz force $Q$ à être constante — impossible car $Q(-\infty)=c\neq0$.
Sinon, $Q$ traverse $x_0$ avec $Q''(x_0)\neq0$, et l'on montre ci-dessous
(Prop.~\ref{prop:fpos}) que ce cas ne se produit pas sur $]0,c[$. Dans tous les
cas $\eps$ est constant.

\subsection{Le second membre ne s'annule pas sur $]0,c[$}
\begin{proposition}\label{prop:fpos}
Pour tout $y\in\,]0,c[$, on a $f(y)>0$.
\end{proposition}

\begin{proof}
Par continuité de $Q$ et puisque $Q(-\infty)=c$, $Q(+\infty)=0$, le théorème des
valeurs intermédiaires assure que $Q$ prend \emph{toute} valeur de $]0,c[$.
Comme $f(Q(x))=Q'(x)^2\ge0$, on a déjà $f\ge0$ sur $]0,c[$.

Supposons par l'absurde $f(y_0)=0$ pour un $y_0\in\,]0,c[$, et soit $x_0$ tel que
$Q(x_0)=y_0$ ; alors $Q'(x_0)=0$. Rappelons que
\[
  f'(y)=2\mu y+2g(y)=2\big(\mu y+g(y)\big),
  \qquad\text{donc}\qquad f'(y_0)=2\,Q''(x_0).
\]
Deux cas se présentent.
\begin{itemize}[leftmargin=1.6em]
  \item \textbf{Si} $\mu y_0+g(y_0)\neq0$ : alors $f'(y_0)\neq0$, donc $f$
  \emph{change de signe} en $y_0$. Il existerait alors des points de $]0,c[$ où
  $f<0$, ce qui contredit $f\ge0$ établi ci-dessus.
  \item \textbf{Si} $\mu y_0+g(y_0)=0$ : alors $y_0$ est un point d'équilibre, et
  la fonction constante $\equiv y_0$ est solution de l'EDO d'ordre 2 avec donnée
  de Cauchy $(y_0,0)$. Par unicité de Cauchy–Lipschitz, $Q\equiv y_0$, ce qui
  contredit $Q(+\infty)=0\neq y_0$.
\end{itemize}
Les deux cas sont absurdes, donc $f(y)>0$ sur $]0,c[$.
\end{proof}

\begin{remarque}
La preuve corrige une erreur fréquente : la fonction constante $y_0$ est solution
de l'EDO \emph{si et seulement si} $\mu y_0+g(y_0)=0$ (équilibre), et \emph{non}
si $f(y_0)=0$. Ces deux conditions ne coïncident qu'aux bords $0$ et $c$,
précisément parce que $f'=2(\mu y+g)$.
\end{remarque}

\subsection{$Q$ est décroissante}
Pour $x$ grand, $Q(x)\to0^+$ et $f(Q(x))>0$ : $Q$ décroît vers $0$ par valeurs
positives, donc $Q'(x)<0$. Comme $\sqrt{f(Q(x))}>0$, ceci impose $\eps=-1$. Ainsi
\[
  Q'(x)=-\sqrt{f(Q(x))}<0\qquad\forall x\in\R,
\]
la stricte négativité venant de la Proposition~\ref{prop:fpos}. La fonction $Q$
est strictement décroissante.

\subsection{Condition d'équilibre en $c$}
En passant à la limite $x\to-\infty$ dans $Q''=\mu Q+g(Q)$ (avec $Q\to c$,
$Q'\to0$, donc $Q''\to0$) :
\[
  0=\mu c+g(c)\quad\Longrightarrow\quad g(c)=-\mu c<0.
\]
Le point $c$ est donc un équilibre de $\mu y+g(y)=0$.

\subsection{Nécessité des conditions de l'algorithme}
On relie maintenant ces résultats à $P$.
\paragraph{$\mu=P(c)$.} En évaluant~\eqref{eq:energie} à la limite $x\to-\infty$ :
$\mu c^2+2G(c)=0$, donc $\mu=-2G(c)/c^2=P(c)$, et $P(c)=\mu>0$.
\paragraph{$P'(c)=0$.} On calcule $P'(x)=-\dfrac{2xg(x)+4G(x)}{x^3}$\ ;
en $x=c$, avec $g(c)=-\mu c$ et $G(c)=-\mu c^2/2$,
\[
  P'(c)=\frac{2\mu c^2-2\mu c^2}{c^3}=0.
\]
\paragraph{$P(x)<P(c)$ sur $]0,c[$.} Pour $x\in\,]0,c[$,
\[
  P(c)-P(x)=\mu+\frac{2G(x)}{x^2}=\frac{\mu x^2+2G(x)}{x^2}=\frac{f(x)}{x^2}>0
\]
par la Proposition~\ref{prop:fpos}. \medskip

\noindent En résumé, toute solution de~\eqref{eq:edo} fournit un couple
$(c,\mu=P(c))$ qui passe les trois tests : \emph{l'algorithme ne manque aucune
solution}.

% =========================================================================
% 4. PARTIE 2 : SYNTHÈSE
% =========================================================================
\section{Partie 2 : Synthèse}
\label{sec:synthese}

On suppose désormais que $c,\mu>0$ vérifient les conditions (i)--(iii), et l'on
\emph{construit} une solution.

\subsection{Positivité de $f$ sous les nouvelles hypothèses}
La condition (iii), $P(x)<P(c)=\mu$ sur $]0,c[$, s'écrit $-2G(x)/x^2<\mu$, soit
\[
  f(x)=\mu x^2+2G(x)>0\qquad\forall x\in\,]0,c[.
\]
La positivité de $f$ est donc une conséquence directe de l'hypothèse (iii).

\subsection{Construction par séparation des variables}
On cherche $Q$ vérifiant $Q'=-\sqrt{f(Q)}$. Pour $y\in\,]0,c[$, $f(y)>0$ donc
$1/\sqrt{f(y)}$ est bien définie. Fixons $y_0\in\,]0,c[$ et posons
\[
  \Phi(y)=\int_{y_0}^{y}\frac{-\dd t}{\sqrt{f(t)}},\qquad y\in\,]0,c[.
\]
$\Phi$ est $\mathcal{C}^1$ et strictement décroissante. Étudions ses limites aux
bords.
\begin{itemize}[leftmargin=1.6em]
  \item En $y\to0^+$ : $f(y)\sim\mu y^2$, donc $1/\sqrt{f(y)}\sim 1/(\sqrt{\mu}\,y)$,
  non intégrable en $0$. Ainsi $\Phi(y)\to+\infty$.
  \item En $y\to c^-$ : on a $f(c)=\mu c^2+2G(c)=0$ (car $\mu=P(c)$) et
  $f'(c)=2(\mu c+g(c))=0$ (condition d'équilibre). Donc $f$ s'annule à l'ordre
  $\ge2$ en $c$, et $1/\sqrt{f}$ n'est pas intégrable en $c$ : $\Phi(y)\to-\infty$.
\end{itemize}
$\Phi:\,]0,c[\,\to\R$ est donc une bijection. On pose $Q=\Phi^{-1}:\R\to\,]0,c[$.
Par construction,
\[
  Q'(x)=-\sqrt{f(Q(x))},\qquad Q(-\infty)=c,\qquad Q(+\infty)=0.
\]

\subsection{$Q$ résout l'équation initiale}
De $Q'^2=f(Q)$, on dérive : $2Q'Q''=f'(Q)Q'=(2\mu Q+2g(Q))Q'$. Comme
$Q'=-\sqrt{f(Q)}\neq0$ sur $\R$, on simplifie :
\[
  Q''=\mu Q+g(Q)\quad\Longleftrightarrow\quad -Q''+\mu Q+g(Q)=0.
\]
Donc $Q$ est bien solution de~\eqref{eq:edo}.

\begin{remarque}
En combinant les Parties 1 et 2, on obtient l'équivalence complète : l'algorithme
est \emph{nécessaire et suffisant}.
\end{remarque}

\subsection{Unicité à translation près}
Soit $\widetilde Q$ une autre solution. Par l'analyse de la Partie 1,
$\widetilde Q'=-\sqrt{f(\widetilde Q)}$. Puisque $Q:\R\to\,]0,c[$ est bijective,
il existe un unique $a\in\R$ tel que $Q(a)=\widetilde Q(0)$. Les fonctions
$x\mapsto\widetilde Q(x)$ et $x\mapsto Q(x+a)$ vérifient la même EDO d'ordre 1
avec la même condition initiale en $0$ ; comme $u\mapsto-\sqrt{f(u)}$ est
$\mathcal{C}^1$ sur $]0,c[$, Cauchy–Lipschitz donne
$\widetilde Q(x)=Q(x+a)$ pour tout $x$.

\begin{theoreme}[Résultat principal]
Le problème~\eqref{eq:edo} admet une solution si et seulement s'il existe $c>0$
tel que $P'(c)=0$, $P(c)>0$ et $P(x)<P(c)$ sur $[0,c[$ ; dans ce cas $\mu=P(c)$
et la solution est unique à translation près.
\end{theoreme}

% =========================================================================
% 5. PARTIE 3 : IMPLÉMENTATION
% =========================================================================
\section{Partie 3 : Implémentation numérique}
\label{sec:numerique}

\subsection{Objectif}
On implémente deux fonctions Python :
\begin{itemize}[leftmargin=1.6em]
  \item \texttt{find\_c\_values(g)} : renvoie tous les couples $(c,\mu)$
  satisfaisant l'algorithme pour une fonction $g$ donnée ;
  \item \texttt{solve\_Q(g, c, mu)} : intègre numériquement
  $Q'=-\sqrt{\mu Q^2+2G(Q)}$ et trace le profil.
\end{itemize}

\subsection{Choix d'un exemple correct : une non-linéarité bistable}
\label{sec:exemple}
Le choix de $g$ est crucial : toutes les fonctions ne produisent pas de front.
Nous prenons
\begin{equation}\label{eq:gtest}
  g(t)=\tfrac12 t^3-\tfrac32 t^2,
  \qquad\text{soit}\qquad
  \mu y+g(y)=\tfrac12\,y(y-1)(y-2).
\end{equation}
Le terme de réaction est une cubique \emph{bistable} (zéros en $0$, $1$, $2$),
caractéristique de l'équation d'Allen–Cahn. On vérifie $g(0)=g'(0)=0$, et
\[
  G(y)=\tfrac18 y^4-\tfrac12 y^3,\qquad
  P(y)=-\frac{2G(y)}{y^2}=y-\frac{y^2}{4},\qquad P'(y)=1-\frac{y}{2}.
\]
Donc $P'(c)=0\Rightarrow c=2$, et $\mu=P(2)=1>0$ ; comme $P$ est une parabole de
sommet $y=2$, la condition (iii) est satisfaite. L'algorithme renvoie
\emph{exactement} $(c,\mu)=(2,1)$.

\paragraph{Solution exacte.} Pour cette famille, $f(y)=\frac{\mu}{c^2}y^2(c-y)^2$,
d'où $Q'=-\frac{\sqrt{\mu}}{c}Q(c-Q)$ (équation logistique). On intègre par
séparation des variables :
\begin{equation}\label{eq:exacte}
  \boxed{\,Q(x)=\frac{c}{1+e^{\sqrt{\mu}\,(x-x_0)}}
        =\frac{c}{2}\left(1-\tanh\frac{\sqrt{\mu}\,(x-x_0)}{2}\right).}
\end{equation}
Cette expression fermée servira de \emph{vérité terrain} pour valider le solveur.

\subsection{Approche pour \texttt{find\_c\_values}}
On échantillonne $[0,c_{\max}]$, on calcule $P$ par intégration numérique
(\texttt{scipy.integrate.quad}), puis $P'$ par différences finies
(\texttt{numpy.gradient}). Les \emph{maxima locaux intérieurs} de $P$ (passage de
$P'>0$ à $P'<0$) sont les candidats ; pour chacun on vérifie (ii) et (iii). Les
candidats détectés en double (deux points de grille voisins) sont fusionnés.

\subsection{Approche pour \texttt{solve\_Q}}
On intègre $Q'=-\sqrt{f(Q)}$ avec \texttt{solve\_ivp} (RK45) en partant du point
milieu $Q(0)=c/2$, séparément vers la droite ($Q\to0$) et vers la gauche
($Q\to c$). Comme $f$ admet un \emph{zéro double} à chaque bord, les paliers sont
des asymptotes : aucune explosion numérique, contrairement à un mauvais choix de
paramètres.

\subsection{Code annoté (cœur algorithmique)}
\noindent\textbf{Détection des paramètres :}
\lstinputlisting[linerange={72-120}]{edo_solver.py}

\noindent\textbf{Intégration du profil :}
\lstinputlisting[linerange={126-149}]{edo_solver.py}

\noindent Le script complet (fonctions $g$ et génération des figures comprises)
est fourni dans le dépôt sous \texttt{edo\_solver.py}.

\subsection{Résultats et illustrations}

\paragraph{Détection des paramètres.}
La figure~\ref{fig:P} montre $P$ pour l'exemple~\eqref{eq:gtest} : un unique
maximum intérieur en $c=2$, donnant $\mu=1$. La zone ombrée illustre la
condition (iii).

\begin{figure}[h]
  \centering
  \includegraphics[width=0.82\linewidth]{fig1_P_bistable.png}
  \caption{Fonction $P$ pour $g(t)=\frac12 t^3-\frac32 t^2$. Le maximum intérieur
  $c=2$ fournit $\mu=P(c)=1$. La zone verte matérialise $P(x)<P(c)$ sur $[0,c[$.}
  \label{fig:P}
\end{figure}

\paragraph{Validation contre la solution exacte.}
La figure~\ref{fig:exact} superpose la solution numérique et la formule
exacte~\eqref{eq:exacte}. L'écart maximal mesuré est de l'ordre de
$7\times10^{-10}$, ce qui valide l'implémentation.

\begin{figure}[h]
  \centering
  \includegraphics[width=0.82\linewidth]{fig3_profil_vs_exact.png}
  \caption{Profil numérique (RK45) \emph{vs.} solution exacte
  $Q(x)=c/(1+e^{\sqrt{\mu}x})$ pour $(c,\mu)=(2,1)$. Écart maximal
  $\approx 7\times10^{-10}$.}
  \label{fig:exact}
\end{figure}

\paragraph{Profils pour deux non-linéarités.}
La figure~\ref{fig:profils} donne les profils pour $(c,\mu)=(2,1)$ et
$(c,\mu)=(3,4.5)$ (famille générale $R(y)=y(y-c/2)(y-c)$). Les deux fronts sont
décroissants et atteignent leurs paliers.

\begin{figure}[h]
  \centering
  \includegraphics[width=0.82\linewidth]{fig4_deux_profils.png}
  \caption{Profils de transition pour deux non-linéarités bistables, avec les
  valeurs de $c$ détectées automatiquement.}
  \label{fig:profils}
\end{figure}

\paragraph{Lecture en plan de phase.}
La figure~\ref{fig:phase} représente l'orbite hétérocline dans le plan $(Q,Q')$ :
elle relie les deux équilibres $0$ et $c$, conformément à $Q'^2=f(Q)$.

\begin{figure}[h]
  \centering
  \includegraphics[width=0.84\linewidth]{fig6_phase.png}
  \caption{Plan de phase : l'orbite (en rouge) connecte les équilibres $0$ et $c$.
  Le profil $Q$ correspond au parcours de cette trajectoire.}
  \label{fig:phase}
\end{figure}

\paragraph{Invariance par translation.}
La figure~\ref{fig:transl} illustre l'unicité à translation : toute translatée
horizontale du front reste solution.

\begin{figure}[h]
  \centering
  \includegraphics[width=0.84\linewidth]{fig5_translation.png}
  \caption{Famille de solutions $Q(\cdot+a)$ pour $a\in\{-5,0,5\}$.}
  \label{fig:transl}
\end{figure}

% =========================================================================
% 6. REGARD CRITIQUE
% =========================================================================
\section{Regard critique et limites}

\paragraph{Quand l'algorithme ne renvoie rien.}
Un point essentiel, souvent négligé, est qu'\emph{une fonction $g$ arbitraire ne
donne pas de front}. Le contre-exemple $g(t)=-t^3$ l'illustre : alors
$G(y)=-y^4/4$ et $P(x)=x^2/2$, strictement croissante, donc $P'(c)=0$ n'a
\emph{aucune} solution $c>0$ (figure~\ref{fig:echec}). L'algorithme renvoie
correctement la liste vide, et de fait le problème n'admet pas de solution :
en plan de phase, l'origine n'est pas connectée à un second équilibre. Ce
comportement, loin d'être un défaut, confirme la \emph{nécessité} démontrée en
Partie 1.

\begin{figure}[h]
  \centering
  \includegraphics[width=0.84\linewidth]{fig2_P_echec.png}
  \caption{Pour $g(t)=-t^3$, $P$ est monotone : pas de maximum intérieur, donc
  aucun couple $(c,\mu)$ admissible. L'algorithme renvoie une liste vide.}
  \label{fig:echec}
\end{figure}

\paragraph{Limites de la méthode numérique.}
\begin{itemize}[leftmargin=1.6em]
  \item La détection des maxima de $P$ repose sur une grille : un maximum très
  pointu ou hors de $[0,c_{\max}]$ peut être manqué. Le paramètre $c_{\max}$ doit
  être choisi avec discernement.
  \item L'intégration de $Q'=-\sqrt{f(Q)}$ près des paliers est sensible car
  $f\to0$ ; on s'appuie sur le caractère asymptotique (zéro double) pour éviter
  les instabilités, mais une tolérance trop lâche dégraderait la précision.
  \item L'estimation de $P'$ par différences finies est bruitée ; une approche
  plus robuste utiliserait une dérivation analytique de $P$ lorsque $g$ est connue
  explicitement.
\end{itemize}

\paragraph{Pistes d'approfondissement.} On pourrait étudier le cas
\emph{non symétrique} (puits de profondeurs différentes), qui fait apparaître une
vitesse de propagation non nulle, ou balayer systématiquement une famille
paramétrée de fonctions $g$ pour cartographier l'existence des fronts.

\subsection{Étude de convergence}
Pour quantifier la fiabilité de la détection, on mesure l'erreur entre le
$c$ renvoyé par \texttt{find\_c\_values} et la valeur exacte $c=2$, en fonction
de la finesse de la grille $n_{\text{grid}}$ (Table~\ref{tab:conv}). L'erreur
décroît régulièrement, approximativement comme $1/n_{\text{grid}}$, ce qui est
cohérent avec l'estimation de $P'$ par différences finies. Par ailleurs, la
solution de \texttt{solve\_Q} reste précise à $\sim10^{-9}$ près quel que soit le
pas, la tolérance \texttt{rtol} du solveur étant le facteur limitant.

\begin{table}[h]
  \centering
  \renewcommand{\arraystretch}{1.2}
  \begin{tabular}{rcc}
    \toprule
    $n_{\text{grid}}$ & $c$ détecté & erreur $|c-2|$\\
    \midrule
    200   & 1.98553 & $1.5\times10^{-2}$\\
    500   & 1.99459 & $5.4\times10^{-3}$\\
    1000  & 1.99760 & $2.4\times10^{-3}$\\
    2000  & 1.99910 & $9.0\times10^{-4}$\\
    4000  & 1.99985 & $1.5\times10^{-4}$\\
    \bottomrule
  \end{tabular}
  \caption{Convergence de la détection de $c$ vers la valeur exacte $2$
  lorsqu'on raffine la grille. L'erreur est divisée par $\sim2$ à chaque
  doublement de $n_{\text{grid}}$.}
  \label{tab:conv}
\end{table}

% =========================================================================
% 6 bis. APPROFONDISSEMENT
% =========================================================================
\section{Approfondissement : rupture de symétrie et vitesse de front}
\label{sec:nonsym}

\subsection{Pourquoi la symétrie est nécessaire}
Notre problème impose un front \emph{stationnaire} (indépendant du temps). Or,
d'après le paragraphe~\ref{sec:maxwell}, ceci n'est possible que si la condition
de Maxwell $\int_0^c R(y)\dd y=0$ est vérifiée. Pour la famille cubique
$R(y)=\beta\,y(y-a)(y-c)$, un calcul direct donne
\[
  \int_0^c R(y)\dd y=\beta\,c^3\,\frac{2a-c}{6},
\]
qui s'annule \emph{si et seulement si} $a=\dfrac{c}{2}$. C'est exactement le cas
symétrique retenu dans~\eqref{eq:gtest} ($a=1$, $c=2$). Dès que $a\neq c/2$, le
front stationnaire n'existe plus : l'algorithme renvoie alors une liste vide,
ce qui est cohérent avec la théorie.

\subsection{Ce que devient le front quand $a\neq c/2$}
Lorsque les deux puits n'ont pas la même profondeur, l'équation d'évolution
$u_t=u_{xx}-R(u)$ admet non pas un front fixe, mais un \emph{front mobile}
$u(t,x)=U(x-\nu t)$ de vitesse $\nu\neq0$. Le profil $U$ vérifie alors
\[
  U''+\nu U'-R(U)=0.
\]
En multipliant par $U'$ et en intégrant sur $\R$ (avec $U'(\pm\infty)=0$), le
terme $\int_\R U'U''$ s'annule et il reste
\[
  \nu\int_\R U'(\xi)^2\dd\xi=-\int_0^c R(y)\dd y.
\]
Comme $\int_\R U'^2>0$, le \emph{signe de la vitesse} $\nu$ est l'opposé de celui
de l'aire $\int_0^c R$ : la phase « énergétiquement la plus profonde » envahit
l'autre. On retrouve $\nu=0$ exactement au point d'équilibre des aires, c'est-à-dire
la condition de Maxwell. Ce résultat, dû à Fife et McLeod~\cite{fifemcleod1977},
montre que notre problème stationnaire est le cas \emph{critique} d'un phénomène
plus général de propagation de fronts.

\subsection{Lien avec la classification des équilibres}
Du point de vue du plan de phase (figure~\ref{fig:phase}), les équilibres $0$ et
$c$ sont des \emph{points selles} de l'EDO d'ordre 2, tandis que l'équilibre
intermédiaire ($y=1$ dans notre exemple) est un \emph{centre}. Le front
stationnaire correspond à l'unique orbite reliant les deux selles : une
\emph{connexion hétérocline}. Cette lecture géométrique rejoint l'analyse de
stabilité des transitions de phase et éclaire pourquoi la solution est unique à
translation près (l'orbite hétérocline est unique, seul le paramétrage temporel
est libre).

% =========================================================================
% 7. CONCLUSION
% =========================================================================
\section{Conclusion}

Ce projet a établi de manière complète l'existence et l'unicité (à translation
près) des solutions de l'EDO~\eqref{eq:edo}, et caractérisé les paramètres
admissibles via la fonction $P$. Les outils clés ont été la \emph{conservation de
l'énergie} (réduction à l'ordre 1), le \emph{théorème de Cauchy–Lipschitz}
(unicité et non-annulation de $f$) et la \emph{séparation des variables}
(construction explicite). L'implémentation Python, \emph{validée contre une
solution exacte} au niveau $10^{-10}$, confirme la théorie et visualise les
profils. Enfin, le rattachement à l'équation d'Allen–Cahn place ce travail dans
le cadre plus large des fronts de réaction–diffusion.

\paragraph{Organisation du travail.}
Le projet a été mené en équipe de cinq. La partie théorique (Parties 1 et 2) a
fait l'objet d'un travail collectif, chaque membre rédigeant et vérifiant un
sous-ensemble des propositions ; la preuve corrigée de la
Proposition~\ref{prop:fpos} a notamment été reprise ensemble après une première
version erronée. L'implémentation Python (Partie 3), la validation contre la
solution exacte et la production des figures ont été réalisées en binôme, le
reste de l'équipe assurant la relecture du code et la rédaction de l'état de
l'art, de l'approfondissement et de l'annexe. La mise en forme \LaTeX{} et
l'harmonisation finale ont été partagées.

% =========================================================================
% ANNEXE
% =========================================================================
\clearpage
\appendix
\section{Annexe : calcul explicite de la solution exacte}
\label{sec:annexe}

Dans le cas bistable~\eqref{eq:gtest}, l'EDO se résout \emph{à la main} et
fournit une solution en forme fermée. C'est cette solution qui sert de
\emph{référence exacte} pour valider l'implémentation numérique
(figure~\ref{fig:exact}). On détaille ici les quatre étapes du calcul, puis on
vérifie que la fonction obtenue est bien solution du problème de départ.

\subsection*{Étape 1 — Forme factorisée de $f$}
On part de $g(t)=\dfrac12 t^3-\dfrac32 t^2$, avec les paramètres détectés
$\mu=1$ et $c=2$. La primitive est
\[
  G(y)=\int_0^y\!\Big(\tfrac12 t^3-\tfrac32 t^2\Big)\dd t
       =\frac{1}{8}y^4-\frac{1}{2}y^3,
\]
d'où
\[
  f(y)=\mu y^2+2G(y)=y^2-y^3+\frac{1}{4}y^4
      =\frac{1}{4}\,y^2\,(y-2)^2
      =\frac{\mu}{c^2}\,y^2\,(c-y)^2.
\]
On retrouve la structure annoncée : $f$ possède un \textbf{zéro double} en $0$
et en $c$, et reste strictement positive entre les deux. Sur $]0,c[$,
\[
  \sqrt{f(y)}=\frac{\sqrt{\mu}}{c}\,y\,(c-y).
\]

\subsection*{Étape 2 — Réduction à une équation à variables séparables}
L'EDO d'ordre 1 issue de l'énergie, $Q'=-\sqrt{f(Q)}$, s'écrit
\[
  \frac{\dd Q}{\dd x}=-\frac{\sqrt{\mu}}{c}\,Q\,(c-Q),
\]
soit, en séparant les variables (licite car $Q\in\,]0,c[$ donne $Q(c-Q)>0$),
\[
  \frac{\dd Q}{Q\,(c-Q)}=-\frac{\sqrt{\mu}}{c}\,\dd x.
\]

\subsection*{Étape 3 — Décomposition en éléments simples et intégration}
On décompose la fraction de gauche :
\[
  \frac{1}{Q\,(c-Q)}=\frac{1}{c}\left(\frac{1}{Q}+\frac{1}{c-Q}\right).
\]
En intégrant membre à membre,
\[
  \frac{1}{c}\Big(\ln Q-\ln(c-Q)\Big)
  =-\frac{\sqrt{\mu}}{c}\,x+\text{cste},
\]
c'est-à-dire, après multiplication par $c$ et en notant $x_0$ la constante,
\[
  \ln\!\frac{Q}{c-Q}=-\sqrt{\mu}\,(x-x_0).
\]

\subsection*{Étape 4 — Inversion}
En exponentiant puis en isolant $Q$ :
\[
  \frac{Q}{c-Q}=e^{-\sqrt{\mu}\,(x-x_0)}
  \;\Longrightarrow\;
  Q=\big(c-Q\big)\,e^{-\sqrt{\mu}\,(x-x_0)}
  \;\Longrightarrow\;
  \boxed{\,Q(x)=\frac{c}{1+e^{\sqrt{\mu}\,(x-x_0)}}
        =\frac{c}{2}\Big(1-\tanh\frac{\sqrt{\mu}}{2}(x-x_0)\Big).}
\]
La constante $x_0$ paramètre la position du front : c'est exactement
l'\emph{invariance par translation} démontrée au paragraphe~3.4.

\subsection*{Vérification}
On contrôle que $Q$ résout bien le problème initial~\eqref{eq:edo}.
\begin{itemize}[leftmargin=1.6em]
  \item \textbf{Conditions aux limites.} Quand $x\to-\infty$, l'exponentielle
  tend vers $0$ donc $Q\to c$ ; quand $x\to+\infty$, elle tend vers $+\infty$
  donc $Q\to0$. \checkmark
  \item \textbf{Équation différentielle.} Posons $\sigma=\sqrt{\mu}$. En dérivant
  la forme logistique, $Q'=-\dfrac{\sigma}{c}\,Q(c-Q)$, puis
  \[
    Q''=-\frac{\sigma}{c}\,(c-2Q)\,Q'
       =\frac{\sigma^2}{c^2}\,Q\,(c-Q)\,(c-2Q).
  \]
  Avec $\mu=1$, $c=2$ : $Q''=\tfrac14 Q(2-Q)(2-2Q)=\tfrac12 Q(Q-1)(Q-2)$, qui est
  précisément $\mu Q+g(Q)=\tfrac12 Q(Q-1)(Q-2)$. Donc $-Q''+\mu Q+g(Q)=0$.
  \checkmark
\end{itemize}
La fonction~\eqref{eq:exacte} est donc rigoureusement solution, ce qui légitime
son emploi comme vérité terrain pour mesurer l'erreur du solveur.

% =========================================================================
% BIBLIOGRAPHIE
% =========================================================================
\clearpage
\begin{thebibliography}{9}
\addcontentsline{toc}{section}{Références}

\bibitem{allencahn1979}
S.~M. Allen et J.~W. Cahn,
\emph{A microscopic theory for antiphase boundary motion and its application to
antiphase domain coarsening}, Acta Metallurgica, \textbf{27}(6):1085--1095, 1979.

\bibitem{fifemcleod1977}
P.~C. Fife et J.~B. McLeod,
\emph{The approach of solutions of nonlinear diffusion equations to travelling
front solutions}, Archive for Rational Mechanics and Analysis,
\textbf{65}(4):335--361, 1977.

\bibitem{aronsonweinberger1978}
D.~G. Aronson et H.~F. Weinberger,
\emph{Multidimensional nonlinear diffusion arising in population genetics},
Advances in Mathematics, \textbf{30}(1):33--76, 1978.

\bibitem{kpp1937}
A.~Kolmogorov, I.~Petrovsky et N.~Piskunov,
\emph{Étude de l'équation de la diffusion avec croissance de la quantité de
matière et son application à un problème biologique}, Bull. Univ. Moscou,
Sér. Int. A, \textbf{1}:1--25, 1937.

\bibitem{nagumo1962}
J.~Nagumo, S.~Arimoto et S.~Yoshizawa,
\emph{An active pulse transmission line simulating nerve axon},
Proceedings of the IRE, \textbf{50}(10):2061--2070, 1962.

\bibitem{evans2010}
L.~C. Evans,
\emph{Partial Differential Equations}, Graduate Studies in Mathematics, vol.~19,
American Mathematical Society, 2\textsuperscript{e} édition, 2010.

\bibitem{strogatz2015}
S.~H. Strogatz,
\emph{Nonlinear Dynamics and Chaos}, Westview Press,
2\textsuperscript{e} édition, 2015.

\end{thebibliography}

\end{document}
